\documentclass[12pt, letterpaper]{article}
\usepackage{geometry}
\usepackage{mhchem}
\usepackage{chemfig}
\begin{document}
\begin{center}
    chemistry honors final exam cheatsheet
\end{center}
\begin{minipage}[t]{0.45\textwidth}
    \begin{tabular}[t]{c|c}
        haloalkane & $\ce{R-X}$ \\
        aldehyde & $\chemfig{R-C(=[1]O)(-[7]H)}$ \\ \\
        ketone & $\chemfig{R-C(=[1]O)(-[7]R)}$ \\
        alcohol & $\ce{R-OH}$ \\
        ether & $\chemfig{R-[1]O-[7]R'}$ \\
        amine & $\ce{R-NH2}$ \\
        carboxylic acid & $\chemfig{R-C(=[1]O)(-[7]OH)}$ \\ \\
        ester & $\chemfig{R-C(=[1]O)(-[7]OR)}$ \\ \\
        amide & $\chemfig{R-C(=[1]O)(-[7]NH2)}$ \\ \\
        acyl halide & $\chemfig{R-C(=[1]O)(-[7]X)}$
    \end{tabular}
\end{minipage}
\begin{minipage}[t]{0.45\textwidth}
    the table to the left is the functional groups. alcohols make things end with -nol. for example: a five-carbon chain with an alcohol group at two is called 2-pentanol. \\
    remember the prefixes: methane, ethane, propane, butane, pentane, hexane, heptane, octane, nonane, decane. ane $\longrightarrow$ ene when it has a double bond somewhere, and $\longrightarrow$ yne when it has a  triple bond somewhere. \\
    acids: if contains hydrogen and another element ends in ate, add hydro-, drop -ide add ic, add -acid to the end. example: $\ce{HNO3}\longrightarrow$ nitric acid. if it ends in -ite, make it -ous. example: $\ce{HNO2}\longrightarrow$nitrous acid.
    strong acids: $\ce{HCl}$, $\ce{HNO3}$, $\ce{H2SO4}$, $\ce{HBr}$, $\ce{HI}$, $\ce{HClO4}$, $\ce{HClO3}$, strong bases: $\ce{LiOH}$, $\ce{NaOH}$, $\ce{KOH}$, $\ce{RbOH}$, $\ce{CsOH}$, $\ce{Ca(OH)2}$, $\ce{Sr(OH)2}$, $\ce{Ba(OH)2}$. PV=nRT, R=0.08206 in atm, 8.314 in kPa, 62.36 in mmHg. equ: $K_{eq}=\frac{[C]^c[D]^D}{[A]^a[B]^b}$, $a$\ce{A} $+$ $b$\ce{B <=>} $c$\ce{C} $+$ $d$\ce{D}, redox reactions give and take electrons, OIL RIG (of electrons) act se: \ce{Li}, \ce{K}, \ce{Ba}, \ce{Ca}, \ce{Na}, \ce{Mg}, \ce{Al}, \ce{Zn}, \ce{Fe}, \ce{Ni}, \ce{Sn}, \ce{Pb}, \ce{H}, \ce{Cu}, \ce{Ag}, \ce{Au}, silver has a charge of +1, zinc has a charge of +2, and cadmium has a charge of +2. determine charges by neutralizing ions. methyl, ethyl, propyl, butyl, pentyl, ...
\end{minipage}


\end{document}